% Exercise 15.5
We must identify the exact solution in order to have something to compare
with our numerical solution. Accordingly, we start by identifying the spring 
force using Hooke's law, so that we have
\[
    F_{\text{spring}}(x_0) = -kx_0,
\]
as in (15.8), with $k = 100,000\,\text{N/m}$. Since we ignore air resistance
and gravity, that is our only force in the differential equation governing
the spring's motion. From Newton's second law, we then have
\[
    F_{\text{spring}}(x) = m\ddot{x},
\]
for $m = 3\,\text{kg}$. Our initial conditions are given by the displacement
from equilibrium, so that $x(0) = 0.01\,\text{m}$ and the ball is released
with no initial velocity, so that $\dot{x(0)} = 0,\text{m/s}$.

Therefore, we have
\[
    m\ddot{x} = -kx,
\]
and since we expect oscillations in the resulting movement, we use the trial
solution $x(t) = e^{rt}$, so that $\ddot{x} = r^2e^{rt}$. Thus we have
\[
    mr^2e^{rt} = -ke^{rt},
\]
and, since $e^{rt} \neq 0$, we have
\begin{equation*}
    \begin{split}
        mr^2 + k &= 0 \\
        \Leftrightarrow r^2 &= -\frac{k}{m} \\
        \Leftrightarrow r &= \pm \sqrt{-\frac{k}{m}} \\
                          &= \pm i \sqrt{\frac{k}{m}},
    \end{split}
\end{equation*}
so that our answer becomes, using Euler's formula and the superposition 
principle,
\[
    x(t) = a\cos{\sqrt{k/m}t} + b\sin{\sqrt{k/m}t}.
\]

Now, applying the first initial condition, we have
\begin{equation*}
    \begin{split}
        x(0) &= a \\
        \Rightarrow a &= 0.01,
    \end{split}
\end{equation*}
and from the second, we have
\begin{equation*}
    \begin{split}
        \dot{x}(0) &= b\sqrt{\frac{k}{m}} \\
        \Rightarrow b &= 0.
    \end{split}
\end{equation*}

Therefore, the exact solution is given by
\begin{equation*}
    \begin{split}
        x(t) &= 0.01\cos{\sqrt{k/m}t} \\
             &= 0.01\cos{\sqrt{\frac{100,000}{3}}t}.
    \end{split}
\end{equation*}

In Figure~\ref{fig15_5}, we have a plot of the mass's position as a function of time,
with labels for the exact solution, as well as the good approximation (with 
$\Delta t = {10}^{-5}\,\text{s}$), and bad approximation (with 
$\Delta t = 5 \times {10}^{-3}\,\text{s}$), using Euler's method.

\begin{figure}[!ht]
    \includegraphics[width=\textwidth, center]{../plots/15_5.pdf}
    \caption{Spring-mass position approximations as a function of time \label{fig15_5}} 
\end{figure}
